% !TeX root = Javier_TARAZONA_CV_fr.tex
% !TeX program = pdflatex
%-------------------------
% Resume in Latex
% Author : Javier Tarazona
% Offre : offre type - Stage Data Scientist (H/F), stage de fin d'études 6 mois (voir offer.md)
%------------------------

\documentclass[a4paper,10pt]{article}
\DeclareUnicodeCharacter{00A0}{~}
\DeclareUnicodeCharacter{2013}{\textendash}
\DeclareUnicodeCharacter{2014}{\textemdash}
\DeclareUnicodeCharacter{201C}{``}
\DeclareUnicodeCharacter{201D}{''}
\DeclareUnicodeCharacter{2022}{\textbullet}
\DeclareUnicodeCharacter{25E6}{\textopenbullet}
\usepackage[T1]{fontenc}      % Codificación de salida (acentos y símbolos)
\usepackage[utf8]{inputenc}   % Lee texto en UTF-8
\usepackage[empty]{fullpage}
\usepackage[pdftex]{hyperref}
\usepackage[usenames,dvipsnames]{color}
\usepackage{array}
\usepackage{enumitem}         % Control de listas
\usepackage{lmodern}
\usepackage{titlesec}
\usepackage{xcolor}

% --- Viñetas limpias y seguras ---
\setlist[itemize]{label=\textbullet, labelsep=0.5em, leftmargin=*}

% ===== CORRECTIFS ATS =====
\input{glyphtounicode}        % chaque glyphe -> vrai texte Unicode (ligatures decomposees)
\pdfgentounicode=1
\usepackage[none]{hyphenat}   % point 5 : plus aucun mot-cle coupe
\sloppy                       %           evite les debordements de marge
\usepackage{textcomp}         % point 6 : apostrophe ASCII au lieu de U+2019
\catcode`\'=\active
\def'{\textquotesingle}
\hypersetup{                  % point 8 : metadonnees indexables
  pdftitle={Javier Andres TARAZONA JIMENEZ - CV - Data Scientist - Stage de fin d'études en Data Science et Machine Learning},
  pdfauthor={Javier Andres TARAZONA JIMENEZ},
  pdfsubject={Data Scientist - stage de fin d'études Bac+5 de 6 mois à partir d'avril 2027 - Data Science, Machine Learning, statistique, Python, SQL},
  pdfkeywords={Data Scientist, Data Science, Machine Learning, apprentissage supervisé, apprentissage non supervisé,
               clustering, modélisation statistique, feature engineering, analyse de données, visualisation de données,
               évaluation de modèles, Deep Learning, NLP, LLM, Python, SQL, pandas, NumPy, scikit-learn, PyTorch,
               Matplotlib, Docker, Git, Google Cloud (GCP), MLOps, probabilités, statistique,
               Bac+5, Diplôme d'Ingénieur, Grande École, stage de fin d'études},
  pdfborder={0 0 0}           % pas de cadre autour des liens (le soulignement bleu suffit)
}
% ==========================

\pagestyle{empty}

% Adjust margins
\addtolength{\oddsidemargin}{-0.45in}
\addtolength{\evensidemargin}{-0.45in}
\addtolength{\textwidth}{1.15in}
\addtolength{\topmargin}{-.62in}
\addtolength{\textheight}{1.20in}

\urlstyle{same}

\raggedbottom
\raggedright

%\definecolor{mydarkgreen}{RGB}{0, 140, 123}
\definecolor{mydarkgreen}{RGB}{0, 0, 0}

% Sections formatting
\titleformat{\section}{
  \vspace{-6pt}\scshape\raggedright\large
}{}{0em}{}[\color{mydarkgreen}\titlerule \vspace{-8pt}]

%-------------------------
% Custom commands

% Indented paragraph helper: \indentbox[ancho]{contenido}
% Uso: \indentbox{Texto...}  o  \indentbox[1.5em]{Texto...}
% [t] + \raggedright + \strut : lignes alignees en haut, sans espaces etires ni chevauchement.
\newcommand{\indentbox}[2][1em]{\hspace*{#1}\parbox[t]{\dimexpr\linewidth-#1\relax}{\raggedright #2\strut}}

% Photo et QR code desactives : dans Workday une photo casse l'analyse du texte voisin.


%-------------------------------------------
%%%%%%  CV STARTS HERE  %%%%%%%%%%%%%%%%%%%%%%%%%%%%


\begin{document}

%----------HEADING-----------------
% Nom seul sur la premiere ligne, puis coordonnees en texte brut dans le corps (pas d'en-tete, pas de minipage).

\noindent{\Large\bfseries\textcolor{mydarkgreen}{Javier Andres TARAZONA JIMENEZ}}\par
\vspace{3pt}
\noindent{\footnotesize javitar06@gmail.com \textbullet{} +33 7 46 36 97 75 \textbullet{} Massy, France}\par
\vspace{1pt}
\noindent{\footnotesize\href{https://www.linkedin.com/in/javier-andres-tarazona-jimenez-84b489222/}{\textcolor{blue}{\underline{linkedin.com/in/javier-andres-tarazona-jimenez-84b489222}}} \textbullet{} \href{https://github.com/JavierTarazona06}{\textcolor{blue}{\underline{github.com/JavierTarazona06}}}}\par

\vspace{0.12cm}

{\centering
  {\large\bfseries Data Scientist - Stage de fin d'études en Data Science et Machine Learning}\\[2pt]
  {\footnotesize Élève ingénieur en dernière année (Grande École, Bac+5) \textbullet{} Disponible à partir d'avril 2027 pour 6 mois}\\[3pt]
  {\footnotesize\textit{Machine Learning, modélisation statistique et analyse de données en Python (pandas, scikit-learn, PyTorch) et SQL}}\par}

\vspace{0.05cm}

%-----------FORMATION-----------------
\section{\textcolor{mydarkgreen}{\textbf{FORMATION}}}
{\raggedright
\textbf{ENSTA Paris (Institut Polytechnique de Paris)} \textbullet{} Palaiseau, France \textbullet{} \textit{\footnotesize 2025/09 - 2027/09}\\
{\footnotesize\textit{Diplôme d'Ingénieur (Bac+5) - Master of Science in Engineering, parcours Intelligence Artificielle}}\\[3pt]
{\footnotesize Cours : Machine Learning, Probability and Machine Learning, Deep Learning, MLOps, Maintenance prédictive, Bases de données}\\[4pt]

\textbf{Universidad Nacional de Colombia (UNAL)} \textbullet{} Bogotá, Colombie \textbullet{} \textit{\footnotesize 2021/01 - 2025/07}\\
{\footnotesize\textit{Cursus d'ingénieur - Ingénierie des Systèmes et Informatique}}\\[3pt]
{\footnotesize 12e université d'Amérique latine (QS 2026). Cours : Probabilités et statistique, Modèles stochastiques, Optimisation, Méthodes numériques}\\[4pt]

\textbf{University of Saskatchewan} \textbullet{} Saskatoon, Canada \textbullet{} \textit{\footnotesize 2024/09 - 2024/12}\\
{\footnotesize\textit{Échange académique - Informatique : Algorithmes, Ingénierie logicielle, Systèmes d'exploitation}}\par}

\vspace{0.08cm}
%-----------EXPERIENCE-----------------
\section{\textcolor{mydarkgreen}{\textbf{EXPÉRIENCE PROFESSIONNELLE}}}
{\raggedright
\textbf{ENSTA Paris - Laboratoire U2IS} \textbullet{} Palaiseau, France \textbullet{} \textit{\footnotesize 2026/05 - 2026/08}\\
{\footnotesize\textit{Stagiaire de recherche en Intelligence Artificielle et Vision 3D}}\\[3pt]
\indentbox{\footnotesize \textbf{- Préparation de données :} pipeline Python/Blender générant des datasets d'entraînement multi-vues (rendu, caméras, métadonnées).}\\
\indentbox{\footnotesize \textbf{- Deep Learning :} pré-entraînement auto-supervisé (\textbf{PyTorch}, \textbf{Vision Transformers}) et curriculum learning sur formes 3D.}\\[4pt]

\textbf{APEDYS91} \textbullet{} Paris, France \textbullet{} \textit{\footnotesize 2025/10 - 2026/02}\\
{\footnotesize\textit{Ingénieur Machine Learning / NLP}}\\[3pt]
\indentbox{\footnotesize \textbf{- NLP et Large Language Models (LLM) :} correction de texte en français par LanguageTool et LLM locaux (Mistral, Gemma).}\\
\indentbox{\footnotesize \textbf{- Fiabilité des modèles :} contrôle des entités et des nombres par reconnaissance d'entités nommées (NER, spaCy) contre les hallucinations.}\\[4pt]

\textbf{Engin A.I} \textbullet{} Bogotá, Colombie (à distance, États-Unis) \textbullet{} \textit{\footnotesize 2025/02 - 2025/06}\\
{\footnotesize\textit{Ingénieur Logiciel et Vision par Ordinateur}}\\[3pt]
\indentbox{\footnotesize \textbf{- Machine Learning :} clustering non supervisé (OpenCV, NumPy) pour la détection de voies : \textbf{+50\% de vitesse, +60\% de précision}.}\\
\indentbox{\footnotesize \textbf{- Analyse de données :} trajectoires de roue (WheelPath) corrélées aux zones de détérioration (+10\% de satisfaction utilisateur).}\\[4pt]

\textbf{Engin A.I} \textbullet{} Bogotá, Colombie (à distance, États-Unis) \textbullet{} \textit{\footnotesize 2024/01 - 2024/07}\\
{\footnotesize\textit{Ingénieur Logiciel et Vision par Ordinateur}}\par}


\vspace{0.08cm}
%-----------PROJETS-----------------
\section{\textcolor{mydarkgreen}{\textbf{PROJETS}}}
{\raggedright
\textbf{Segmentation de peau par apprentissage statistique} \textbullet{} {\footnotesize Équipe de 3, 3 mois, rapport en ligne}\\
{\footnotesize\href{https://github.com/JavierTarazona06/Vision_LocalFeatures_BayesianClassification_KMeans}{\textcolor{blue}{\underline{github.com/JavierTarazona06/Vision\_LocalFeatures\_BayesianClassification\_KMeans}}}}\\
\indentbox{\footnotesize \textbf{- Feature engineering :} variables RGB/HSV/YCrCb et Cb-Cr enrichies par le gradient de luminance pour la classification pixel à pixel.}\\
\indentbox{\footnotesize \textbf{- Modélisation :} comparaison de Gaussian Naive Bayes et QDA (supervisé) avec K-Means (non supervisé) ; scikit-learn, Matplotlib.}\\[3pt]

\textbf{"Tameo" - Vision embarquée pour bateau autonome} \textbullet{} \textit{\footnotesize 2025/09 - 2027/06}\\
{\footnotesize\textit{Équipe de 5, membre puis chef de projet depuis 2026/09 - Monaco Energy Boat Challenge, AI Class - Python, PyTorch, YOLO}}\\
\indentbox{\footnotesize \textbf{- Évaluation de modèles :} fine-tuning et data augmentation de détecteurs ; mAP, F1, précision/rappel portés d'environ 80\% à 90\%.}\\[3pt]

\textbf{"Sharp Sight" - Pipeline de collecte et de traitement de données} \textbullet{} {\footnotesize\href{https://github.com/JavierTarazona06/SharpSight}{\textcolor{blue}{\underline{github.com/JavierTarazona06/SharpSight}}}}\\
\indentbox{\footnotesize \textbf{- Données :} web scraping (Selenium) multi-sites e-commerce ; nettoyage, structuration et catégorisation avec pandas ; API REST FastAPI.}\\[3pt]

\textbf{"ORIUN" - Plateforme web de gestion de candidatures} \textbullet{} {\footnotesize\href{https://github.com/JavierTarazona06/ORIUN_back}{\textcolor{blue}{\underline{github.com/JavierTarazona06/ORIUN\_back}}}}\\
\indentbox{\footnotesize \textbf{- Données et déploiement :} API REST Django, base PostgreSQL (SQL), statistiques d'aide à la décision ; Docker, Google Cloud.}\par}


\vspace{0.08cm}
%-----------COMPETENCES TECHNIQUES-----------------
\section{\textcolor{mydarkgreen}{\textbf{COMPÉTENCES TECHNIQUES}}}
{\footnotesize
\setlength{\parindent}{0pt}\setlength{\parskip}{2pt}
\hangindent=1.4em\hangafter=1 \textbf{Data Science :} Machine Learning (supervisé, non supervisé, clustering), modélisation statistique, feature engineering,\newline analyse et visualisation de données, évaluation de modèles (précision, rappel, F1), Deep Learning, NLP, LLM\par
\hangindent=1.4em\hangafter=1 \textbf{Langages et librairies :} Python, SQL, pandas, NumPy, scikit-learn, PyTorch, Matplotlib, spaCy, OpenCV, C++/C, Java\par
\hangindent=1.4em\hangafter=1 \textbf{Mathématiques :} probabilités et statistique, modèles stochastiques, optimisation, méthodes numériques, algèbre linéaire\par
\hangindent=1.4em\hangafter=1 \textbf{Données et industrialisation :} PostgreSQL, MySQL, Docker, FastAPI, Git/GitHub, Google Cloud (GCP), CI/CD, Agile (SCRUM)\par
}

\vspace{0.08cm}
%-----------LANGUES-----------------
\section{\textcolor{mydarkgreen}{\textbf{LANGUES}}}
{\footnotesize
\noindent \textbf{Anglais :} courant (C1 - IELTS 2024) \textbullet{} \textbf{Français :} avancé (B2 - DELF 2024) \textbullet{} \textbf{Espagnol :} langue maternelle \textbullet{} \textbf{Portugais :} notions
}

\vspace{0.08cm}
%--------DISTINCTIONS------------
\section{\textcolor{mydarkgreen}{\textbf{DISTINCTIONS}}}
{\footnotesize
\noindent \textbf{Bourse Eiffel Excellence} \textbullet{} bourse d'excellence du gouvernement français pour le Master à ENSTA Paris\\[1pt]
\noindent \textbf{Mention d'Honneur - ICPC Colombie 2023} \textbullet{} XXXVII Marathon National de Programmation ACIS REDIS\\[1pt]
\noindent \textbf{Excellence académique} \textbullet{} UNAL, exonération des frais de scolarité
}


%-------------------------------------------
\end{document}
